\documentclass[10pt]{article}
\usepackage[T1]{fontenc}
\usepackage[utf8]{inputenc}
\usepackage[french]{babel}
\usepackage{amsmath,amssymb,mathrsfs,mathpazo}
\usepackage{geometry}
\geometry{a4paper,margin=2cm}
\usepackage{graphicx,array,tabularx,tabularray,colortbl,multirow,multicol,booktabs,enumitem,pifont,ulem,fancybox,diagbox,cancel}
\usepackage{pstricks,pstricks-add,pst-plot,pst-tree,pst-node,pst-text,pst-func,pst-math,pst-bspline,pst-3dplot,pst-eucl,pst-solides3d}
\let\pscalculate\undefined
\usepackage{pst-bezier}
\usepackage{tikz,pgfplots,tkz-tab}
\usetikzlibrary{arrows,arrows.meta,patterns,calc,positioning,trees,decorations.markings,shapes.geometric,shapes.misc,shapes.arrows,intersections,decorations.pathmorphing,decorations.pathreplacing,backgrounds,shadows,babel}
\pgfplotsset{compat=1.18}
\usepgfplotslibrary{fillbetween}
\usepackage[np]{numprint}
\usepackage{listings,makecell,pas-tableur,twoopt}
\let\barre\undefined
\usepackage[pstricks]{tablvar}
\definecolor{darkgreen}{RGB}{0,100,0}
\providecommand{\red}{\color{red}}
\providecommand{\blue}{\color{blue}}
\providecommand{\green}{\color{green}}
\providecommand{\black}{\color{black}}
\providecommand{\white}{\color{white}}
\providecommand{\gray}{\color{gray}}
\providecommand{\cyan}{\color{cyan}}
\providecommand{\magenta}{\color{magenta}}
\providecommand{\yellow}{\color{yellow}}
\providecommand{\violet}{\color{violet}}
\providecommand{\orange}{\color{orange}}
\providecommand{\euro}{\texteuro}
\providecommand{\EUR}{\texteuro}
\providecommand{\ieme}{\textsuperscript{e}}
\providecommand{\degres}{\ensuremath{^\circ}}
\providecommand{\degree}{\ensuremath{^\circ}}
\providecommand{\vv}[1]{\overrightarrow{#1}}
\providecommand{\Coord}[1]{#1}
\providecommand{\decosix}{}
\providecommand{\starredbullet}{\ensuremath{\bullet}}
\pagestyle{empty}
\setlength{\parindent}{0pt}
\usepackage{adjustbox}\DeclareUnicodeCharacter{2236}{:}\DeclareUnicodeCharacter{2190}{\ensuremath{\leftarrow}}
\providecommand{\R}{\mathbb{R}}
\providecommand{\N}{\mathbb{N}}
\providecommand{\Z}{\mathbb{Z}}
\providecommand{\Q}{\mathbb{Q}}
\providecommand{\D}{\mathbb{D}}
\providecommand{\C}{\mathbb{C}}
\providecommand{\e}{{\mathrm{e}}}
\providecommand{\vect}[1]{\overrightarrow{#1}}
\providecommand{\vectt}[1]{\overrightarrow{\mathrm{#1}}}
\providecommand{\vvt}[1]{\overrightarrow{\mathrm{#1}}}
\providecommand{\np}[2][]{{#2~#1}}
\providecommand{\vv}[1]{{\overrightarrow{#1}}}
\renewcommand{\d}{{\mathrm{d}}}
\providecommand{\degres}{{^\circ}}
\providecommand{\degre}{{^\circ}}
\providecommand{\blue}{\color{blue}}
\providecommand{\red}{\color{red}}
\providecommand{\enskip}{\quad}
\providecommand{\llbracket}{\lbrack\!\lbrack}
\providecommand{\rrbracket}{\rbrack\!\rbrack}
\providecommand{\euro}{{\text{€}}}
\providecommand{\no}{{\text{n°}}}
\providecommand{\interval}[2]{{[#1;#2]}}
\providecommand{\intervFO}[2]{{[#1;#2[}}
\providecommand{\dotfill}{\ldots}

\providecommand{\ds}{\displaystyle}
\providecommand{\textding}[1]{#1}
\providecommand{\Oij}{\(\mathrm{O};\vec{\imath},\vec{\jmath}\)}
\providecommand{\Oijk}{\(\mathrm{O};\vec{\imath},\vec{\jmath},\vec{k}\)}
\providecommand{\Ouv}{\(\mathrm{O};\vec{u},\vec{v}\)}
\providecommand{\intervFF}[2]{\ensuremath{[#1;#2]}}
\renewcommand{\euro}{\ifmmode\text{\texteuro}\else\texteuro\fi}
\DeclareUnicodeCharacter{00B0}{\ifmmode^{\circ}\else\textdegree\fi}
\DeclareUnicodeCharacter{20AC}{\ifmmode\text{\texteuro}\else\texteuro\fi}
% Énoncés : Éducation nationale. Transcription et corrigés : APMEP, Denis Vergès et contributeurs. Découpage et classement : CoopMaths, Sébastien Lozano, Éric Elter et contributeurs. Adaptation de la présentation : Maths Explorer.
% CC BY-SA : https://coopmaths.fr/www/about/
% https://coopmaths.fr/alea/static/eam/2026/tex/eam_specialite_2026_06_etrangers_1.tex
\begin{document}
\noindent\textit{APMEP / CoopMaths. CC BY-SA. Présentation adaptée par Maths Explorer.}\par\medskip
Un maire souhaite végétaliser sa ville. Pour cela, il décide de planter des mûriers platanes dans les différents parcs.

Ces arbres sont réputés pour leurs qualités d'ombrage et de résistance à la sécheresse.

\subsection*{Partie A}

Au moment de la plantation, un mûrier platane mesure 1 mètre.

On suppose que chaque année la hauteur de l'arbre augmente de 40 cm.

Pour tout entier naturel $n$, on note $u_n$ la hauteur de l'arbre, en mètres, $n$ années après sa plantation. Ainsi $u_0=1$.

\begin{enumerate}
\item \begin{enumerate}
\item Calculer $u_1$.
\item Quelle sera la hauteur de l'arbre deux années après sa plantation ?
\end{enumerate}
\item Quelle est la nature de la suite $(u_n)$ ? Justifier.
\item Pour tout entier naturel $n$, exprimer $u_n$ en fonction de $n$.
\item Au bout de combien d'années le mûrier atteindra-t-il 9 mètres de haut ?
\end{enumerate}


\subsection*{Partie B}

Au moment de la plantation, l'arbre possède un tronc et deux branches.

Un an après la plantation, on observe 4 nouvelles branches.

Deux ans après la plantation, on observe 8 nouvelles branches.

Chaque année, le nombre de nouvelles branches double comme représenté sur le schéma ci-dessous :

\begin{center}
	\begin{tikzpicture}[
		scale=0.9,
		point/.style={circle, draw=black, fill=gray!65, inner sep=1.8pt},
		branche/.style={line width=1.3pt, gray!75!black},
		sol/.style={line width=1.2pt, teal!70!black}
		]
		
		
		\draw[sol] (-1.7,0) -- (12.7,0);
		
		
		
		
		\coordinate (T0) at (1,0);
		\coordinate (N0) at (1,1.45);
		\coordinate (G0) at (-0.15,2.05);
		\coordinate (D0) at (2.15,2.05);
		
		\draw[branche] (T0) -- (N0);
		\draw[branche] (N0) -- (G0);
		\draw[branche] (N0) -- (D0);
		
		\node[point] at (T0) {};
		\node[point] at (N0) {};
		\node[point] at (G0) {};
		\node[point] at (D0) {};
		
		\node[left] at (-0.35,1.5) {branche};
		\node[right] at (1.7,1.5) {branche};
		\node[left] at (0.65,0.45) {tronc};
		
		\node[align=center] at (1,-1) {au moment\\de la plantation\\(année 0)};
		
		
		
		
		\coordinate (T1) at (5.2,0);
		\coordinate (N1) at (5.2,1.45);
		\coordinate (G1) at (4.05,2.05);
		\coordinate (D1) at (6.35,2.05);
		\coordinate (GG1) at (3.5,3.1);
		\coordinate (GD1) at (4.6,3.1);
		\coordinate (DG1) at (5.8,3.1);
		\coordinate (DD1) at (6.9,3.1);
		
		\draw[branche] (T1) -- (N1);
		\draw[branche] (N1) -- (G1);
		\draw[branche] (N1) -- (D1);
		\draw[branche] (G1) -- (GG1);
		\draw[branche] (G1) -- (GD1);
		\draw[branche] (D1) -- (DG1);
		\draw[branche] (D1) -- (DD1);
		
		\foreach \P in {T1,N1,G1,D1,GG1,GD1,DG1,DD1}
		{
			\node[point] at (\P) {};
		}
		
		\node[align=center] at (5.2,-1) {un an après\\la plantation\\(année 1)};
		
		
		
		
		\coordinate (T2) at (10.1,0);
		\coordinate (N2) at (10.1,1.6);
		
		\coordinate (G2) at (8.85,2.3);
		\coordinate (D2) at (11.35,2.3);
		
		\coordinate (GG2) at (8.2,3.25);
		\coordinate (GD2) at (9.15,3.25);
		\coordinate (DG2) at (10.85,3.25);
		\coordinate (DD2) at (11.8,3.25);
		
		\coordinate (A2) at (7.8,4.2);
		\coordinate (B2) at (8.55,4.2);
		\coordinate (C2) at (8.9,4.2);
		\coordinate (D2top) at (9.65,4.2);
		
		\coordinate (E2) at (10.45,4.2);
		\coordinate (F2) at (11.2,4.2);
		\coordinate (G2top) at (11.55,4.2);
		\coordinate (H2) at (12.3,4.2);
		
		\draw[branche] (T2) -- (N2);
		\draw[branche] (N2) -- (G2);
		\draw[branche] (N2) -- (D2);
		
		\draw[branche] (G2) -- (GG2);
		\draw[branche] (G2) -- (GD2);
		\draw[branche] (D2) -- (DG2);
		\draw[branche] (D2) -- (DD2);
		
		\draw[branche] (GG2) -- (A2);
		\draw[branche] (GG2) -- (B2);
		\draw[branche] (GD2) -- (C2);
		\draw[branche] (GD2) -- (D2top);
		
		\draw[branche] (DG2) -- (E2);
		\draw[branche] (DG2) -- (F2);
		\draw[branche] (DD2) -- (G2top);
		\draw[branche] (DD2) -- (H2);
		
		\foreach \P in {T2,N2,G2,D2,GG2,GD2,DG2,DD2,A2,B2,C2,D2top,E2,F2,G2top,H2}
		{
			\node[point] at (\P) {};
		}
		
		\node[align=center] at (10.1,-1) {deux ans après\\la plantation\\(année 2)};
		
	\end{tikzpicture}
\end{center}

Pour tout entier naturel $n$, on note $v_n$ le nombre de nouvelles branches $n$ années après la plantation. À la plantation, l'arbre possède 2 branches, ainsi on pose $v_0=2$.

\begin{enumerate}
\item Quelle est la nature de la suite $(v_n)$ ? Justifier.
\item \begin{enumerate}
\item Un an après la plantation, l'arbre a produit 4 nouvelles branches. Il possède alors un nombre total de branches égal à 6.

Montrer que, trois ans après sa plantation, l'arbre possède un nombre total de branches égal à 30.
\item On donne le programme ci-contre écrit en langage Python :

\begin{minipage}[t]{0.48\linewidth}
\emph{Rappel :} \texttt{for i in range(n)} permet de répéter $n$ fois un ensemble d'instructions.
\end{minipage}\hfill
\begin{minipage}[t]{0.42\linewidth}
\fbox{\begin{minipage}{0.9\linewidth}\ttfamily
v = 2\par
total = 2\par
for i in range(10):\par
\quad v = 2 * v\par
\quad total = total + v\par
print(total)
\end{minipage}}
\end{minipage}

La valeur affichée par ce programme est 4 094.

Dans le contexte de l'exercice, que représente la valeur 4 094 affichée par ce programme ?
\end{enumerate}
\end{enumerate}

\clearpage

\end{document}
